% ============================================================================
\chapter{Conclusion}\label{ch:conclusion}
% ============================================================================

This thesis set out to benchmark three families of implied-volatility smoothers, and ended by
changing the question the benchmark asks.

The parametric study came first. Implementing Gatheral and Jacquier's theory and reproducing
their published values gave us both a validated toolkit and a reference range: per-slice SVI
fits held-out quotes to $0.64\,\vp$ but leaks butterfly arbitrage on $12\%$ of slices, while
certified SSVI is arbitrage-free by Theorem~\ref{lem:range} and pays roughly three times the
error for it. Any learning-based method therefore needs to operate within this range to offer a
meaningful improvement over the parametric baselines.

The deep smoother landed inside it, and produced the observation around which the rest of the
thesis was organised. On a stress designed to force a calendar violation, the calendar penalty
read exactly zero at its collocation nodes while an independent audit found violations on $28\%$
of the domain. Both numbers were correct. The cause was the grid: exponentially spaced
maturities, which the literature recommends because the butterfly constraint bites at the short
end, leave a gap of nearly half a year at the long end, and the optimiser satisfied the penalty
at every node while letting the surface dip between them. Theorem~\ref{prop:nodegap} formalises
this mechanism by bounding the size of a violation that can remain hidden between neighbouring
nodes. Serving each constraint from its own grid reduced material violations from $34.4\%$ to
zero on the same stress.

The operator study then showed that the problem is not confined to one architecture or to one
kind of grid. A soft penalty only controls the locations at which it is evaluated. For an
operator, that limitation extends across days as well, since satisfying the penalty on training
surfaces does not guarantee the same behaviour on unseen ones. The most accurate model in the
study, the GNO, is materially butterfly-violating on $355$ of $386$ out-of-sample days while its
own training-grid audit reads clean.

Chapter~\ref{ch:downstream} asked what that costs. Through
$\sigma_{\mathrm{loc}}^2=\partial_\tau w/g$, these violations translate directly into
local-volatility defects: roughly half of the GNO surface requires repair, and the associated
risk-neutral densities carry up to $230\%$ negative mass. The same deterioration is visible in
the residual study, where the clean surface outperforms its arbitrage-dirty counterpart in every
matched backtest cell. A separate
experiment showed that the market's own spread structure prices out residual mispricings below
about three volatility points, so the value of a clean surface is in marking and market-making
rather than in taker-side trading.

Taken together, these results expose a clear trade-off: the model with the best vanilla fit is
also the one with the weakest downstream validity. The prior-embedded operator of
Chapter~\ref{ch:methods} addresses it. Fitting a certified SSVI prior per day, and letting the
operator learn only a bounded multiplicative correction, removes every material violation across
all $386$ test days while improving the fit from $2.69$ to $1.32\,\vp$. In this experiment the
added protection therefore comes without a loss of fit.

\paragraph{Main implications.} Three points seem to us to generalise beyond this data set.
First, a soft no-arbitrage constraint is only as informative as the set on which it was
evaluated, so independent auditing must extend beyond both the collocation nodes and the
training dates. Second, where a closed-form certificate exists, incorporating it directly into
the model provides stronger protection than approximating it with a penalty. Finally, the audit
should be evaluated in downstream pricing quantities rather than only through violation counts,
because the same derivatives controlled by the no-arbitrage conditions enter directly into
densities and local volatility.

The economic reading is more specific than we expected at the outset. Arbitrage violations in a
fitted surface are not a trading opportunity the market failed to remove: our scanner found
essentially no executable static arbitrage in eight years of SPX quotes. They are better
understood as a source of model risk, one that becomes most visible when the surface is used for
path-dependent products, because such payoffs depend on the surface in regions that vanilla
quotes constrain only weakly.

\paragraph{Future work.} The results suggest three directions for further work.

The first priority is to complete and verify the experiments that are already in place.
Notebook~06 contains a path-dependent experiment: a down-and-out call priced under an
independent Heston ground truth, with a finite-difference cross-check and a paired hedging
simulation. Its preliminary results are consistent with the findings reported above and suggest
that the downstream effects may be larger for path-dependent products. Its numerical components
could not be verified to the standard applied to the rest of this thesis in the time available,
so we have not reported them as findings, and completing that verification is the immediate next
step.

The same applies to several larger experiments that were begun but not finished. These include
the full multi-year production run, which is intended to replace the staging tables and to test the
``butterfly binds in stress'' hypothesis on the Volmageddon and COVID weeks, the full
residual-economics sample, and the VIX variance-swap replication started in Notebook~05. The
last of these is particularly useful, since it would provide an external benchmark for the
behaviour of the fitted wings.

The second direction is to test the framework in other markets. Much of the pipeline, including
the cleaning, the independent audit and the certified-prior architecture, can be carried over to
bitcoin options. That market provides a useful contrast with SPX: quotes are generally sparser,
spreads wider, maturities shorter, and static-arbitrage violations potentially more prevalent.
It may also change the economics of the frontier result, since wider spreads raise the breakeven
while larger and more persistent mispricings may clear it. A particularly interesting question
is whether the prior-embedded architecture retains the same protection where SSVI provides a
weaker description of the underlying smile.

The third direction concerns new model architectures. Several extensions suggested by
Dr~Arthur Bizzi remain to be explored. One possibility is to transport operator constraints
across days, either by penalising on generated query days or by introducing certified projection
layers. Another is to complete the prior-embedded family with the graph operator, whose training
run could not be retained in the present study. A third is to smooth local volatility directly,
where Proposition~\ref{prop:dupire} makes the constraint set explicit and
Theorem~\ref{prop:nodegap} provides a natural starting point for adaptive collocation with a
provable node-versus-audit bound. The recent non-parametric SANOS framework \cite{sanos2026}
offers another natural setting in which to test the independent-audit methodology developed
here.